Boreal Network serves as the concrete artifact for this paper. It is not presented as a finished production system that already proves every downstream claim. Rather, it provides a structured and inspectable baseline for studying request-rooted orchestration. The artifact is docs-first: canonical object meanings, lifecycle rules, event semantics, request-processing rules, and governance constraints are defined in root documentation before or alongside implementation work \cite{boreal_repo_2026,boreal_network_thesis_2026,boreal_request_processing_2026}.

The artifact also includes machine-readable layers. Canonical object shapes are represented through JSON Schema. HTTP surfaces are represented through OpenAPI. Durable request activity and related event channels are represented through event contracts. Deterministic fixtures provide canonical request threads and planner or matcher scenarios. An eval runner checks extraction, routing, matching, planning, and mutation-safety expectations against those fixtures. Together, these layers make the system more than a prose concept. They provide a reproducible contract baseline that can be inspected, extended, and tested.

This structure is helpful for research because it separates semantic claims from implementation-local behavior. The object model and invariants can be discussed directly. The planner rules can be evaluated against fixed scenarios. The lifecycle and mutation boundaries can be tested without requiring that every future product surface or peer runtime already be fully built. In that sense, Boreal is a suitable artifact for studying how a request-rooted model behaves even before large-scale deployment data exists.

For the current paper, the artifact should be read as a prototype baseline plus evaluation substrate. The research claims are therefore bounded: Boreal demonstrates a coherent durable model, a contract stack, and an evaluation direction. It does not, by itself, establish market-wide superiority, complete empirical validation, or coverage of every class of embodied work. Those stronger claims belong to future studies.
